\documentclass[11pt]{article}
\usepackage[margin=1in]{geometry}
\usepackage{amsmath,amssymb}
\usepackage{enumitem}
\newcommand{\E}{\mathrm{E}}
\newcommand{\Var}{\mathrm{Var}}
\usepackage{graphicx}
\usepackage{booktabs}
\usepackage{bookmark}
\hypersetup{
  pdftitle={PSET 9: Classical statistical inference},
  hidelinks}

\title{PSET 9: Classical statistical inference}
\author{}
\date{\vspace{-2.5em}}

% Problems are numbered 1-8 across the document.
\newcounter{problem}
\newcommand{\problem}[1]{\refstepcounter{problem}\section*{Problem \theproblem: #1}}

\begin{document}
\maketitle

% Shared assignment questions (header block).
% Include at the top of any problem set with:  % Shared assignment questions (header block).
% Include at the top of any problem set with:  % Shared assignment questions (header block).
% Include at the top of any problem set with:  \input{assignment-qs}
\section*{Assignment Qs}

\begin{itemize}
\item Name
\item How long did this problem set take you (round to nearest hour)? \quad
  0-1 hour \ 2-3 hours \ 4-5 hours \ 5+ hours
\item How difficult was this problem set? \quad
  very easy \ 1 \ 2 \ 3 \ 4 \ 5 \ very challenging
\end{itemize}

\section*{Assignment Qs}

\begin{itemize}
\item Name
\item How long did this problem set take you (round to nearest hour)? \quad
  0-1 hour \ 2-3 hours \ 4-5 hours \ 5+ hours
\item How difficult was this problem set? \quad
  very easy \ 1 \ 2 \ 3 \ 4 \ 5 \ very challenging
\end{itemize}

\section*{Assignment Qs}

\begin{itemize}
\item Name
\item How long did this problem set take you (round to nearest hour)? \quad
  0-1 hour \ 2-3 hours \ 4-5 hours \ 5+ hours
\item How difficult was this problem set? \quad
  very easy \ 1 \ 2 \ 3 \ 4 \ 5 \ very challenging
\end{itemize}


\problem{Properties of estimators}

\begin{enumerate}[label=\alph*.]
\item Let \(X_1, \ldots, X_n \sim \text{Poisson}(\lambda)\) and let
  \(\hat{\lambda} = \dfrac{\sum_{i=1}^n X_i}{2n}\). Find the bias,
  standard error, and MSE of this estimator.\footnote{Inspired by
    Wasserman 6.6.1}
\item Let \(X_1, \ldots, X_n \sim \text{Uniform}(0, \theta)\) and let
  \(\hat{\theta} = 3 \bar{X}_n\). Find the bias, standard error, and
  MSE of this estimator. (Recall that if \(X \sim \text{Uniform}(0,
  \theta)\), then \(\E[X] = \theta/2\) and \(\Var(X) =
  \theta^2/12\).)\footnote{Inspired by Wasserman 6.6.3}
\end{enumerate}

\problem{Social newsing}\label{prob:socialnews}

A poll conducted in 2025 found that 12\% of U.S. adults get at least
some news on TikTok. The standard error for this estimate was 2.4\%,
and a normal distribution may be used to model the sample
proportion.\footnote{Inspired by OI 4.8 and 4.10}

Construct a 95\% confidence interval for the fraction of U.S. adults
who get some news on TikTok, and interpret the confidence interval in
context.

\problem{(T/F) Social newsing}

Identify the following statements as true or false using the
information from Problem~\ref{prob:socialnews}. Provide an explanation
to justify each of your answers. Explain in approximately 75 words.

\begin{enumerate}[label=\alph*.]
\item The data provide statistically significant evidence that more
  than 16\% of U.S. adults get some news through TikTok. Use a
  significance level of \(\alpha = 0.05\).
\item Since the standard error is 2.4\%, we can conclude that 97.6\%
  of all U.S. adults were included in the study.
\item If we want to reduce the standard error of the estimate, we
  should collect less data.
\end{enumerate}

\problem{Dating on college campuses}

A survey conducted on a reasonably random sample of 200 undergraduates
asked, among many other questions, about the number of exclusive
relationships these students have been in. The sample average is
3.2 with a standard deviation of 2.17.\footnote{Inspired by OI 4.15}

Estimate the average number of exclusive relationships undergraduate
students have been in using the Normal distribution and a 99\%
confidence interval, and interpret this interval in context. Would a
95\% interval be wider or narrower than the one you found? Explain in
one sentence.

\problem{Statistical significance}

Determine whether the following statement is true or false, and
explain your reasoning: ``With large sample sizes, even small
differences between the null value and the point estimate can be
statistically significant.''\footnote{Inspired by OI 4.47}

\problem{Sleep deprivation}

New York is known as ``the city that never sleeps''. Eight hours is
the standard nightly recommendation. A random sample of 25 New Yorkers
were asked how much sleep they get per night; statistical summaries of
these data are shown below. Do these data provide strong evidence that
the average amount New Yorkers sleep differs from 8 hours a night? Use
a \emph{two-sided} test with \(\alpha = 0.05\), per our course
convention.\footnote{Inspired by OI 5.7}

\begin{center}
\begin{tabular}{ccccc}
\toprule
\(n\) & \(\bar{x}\) & \(s\) & \(\min\) & \(\max\) \\
\midrule
\(25\) & \(7.63\) & \(0.87\) & \(6.17\) & \(9.78\) \\
\bottomrule
\end{tabular}
\end{center}

Critical values of the \(t\) distribution with 24 degrees of freedom:

\begin{center}
\begin{tabular}{lcccc}
\toprule
two-tailed \(\alpha\) & 0.10 & 0.05 & 0.02 & 0.01 \\
\midrule
\(t^*\) & 1.711 & 2.064 & 2.492 & 2.797 \\
\bottomrule
\end{tabular}
\end{center}

\begin{enumerate}[label=\alph*.]
\item Write the hypotheses in symbols and in words.
\item Calculate the test statistic, \(T\), and state the degrees of
  freedom.
\item Using the table above, bound the p-value between two of the
  listed \(\alpha\) levels and interpret it in this context.
\item What is the conclusion of the hypothesis test?
\item Suppose that New Yorkers really do sleep 8 hours a night on
  average. What type of error would your conclusion in part (d)
  represent? Describe what the other type of error would look like in
  this setting.
\end{enumerate}

\problem{Interpreting public opinion polls}

On June 28, 2012 the U.S. Supreme Court upheld the much debated 2010
healthcare law, declaring it constitutional. A Gallup poll released
the day after this decision indicates that 46\% of 1,012 Americans
agree with this decision. At a 95\% confidence level, this sample has
a 3\% margin of error. Based on this information, determine if the
following statements are true or false, and explain your
reasoning.\footnote{Inspired by OI 6.6}

\begin{enumerate}[label=\alph*.]
\item We are 95\% confident that between 43\% and 49\% of Americans in
  this sample support the decision of the U.S. Supreme Court on the
  2010 healthcare law.
\item We are 95\% confident that between 43\% and 49\% of Americans
  support the decision of the U.S. Supreme Court on the 2010
  healthcare law.
\item If we considered many random samples of 1,012 Americans, and we
  calculated the sample proportions of those who support the decision
  of the U.S. Supreme Court, 95\% of those sample proportions will be
  between 43\% and 49\%.
\item The margin of error at a 90\% confidence level would be higher
  than 3\%.
\end{enumerate}

\problem{Power and sample size}

Suppose you have two samples: one with a sample mean of 12, a standard
deviation of 2.13, and \(n_1\) of 87; and a second with a sample mean
of 12.8, a standard deviation of 3.25, and \(n_2\) of 89. Throughout,
use a two-sided \emph{Wald test} of \(H_0: \delta = 0\) versus
\(H_1: \delta \neq 0\), where \(\delta = \mu_1 - \mu_2\).

\begin{enumerate}[label=\alph*.]
\item Is it statistically significant?
\item Now suppose the true difference in means is in fact 0.8. Using
  the power function for the Wald test from class, find the power of
  this test at \(\alpha = 0.05\). What does your answer say about how
  much to trust your conclusion in part (a)?
\item What is the smallest \(n_2\) you could have if you want to
  reject the null hypothesis at \(\alpha < 0.05\)?
\item Darn it! Your \(n_2\) is actually only 85! What should your goal
  be: smaller sample deviation by 10\% or larger sample by 10? Explain
  your reasoning and provide mathematical support.
\end{enumerate}

% Shared AI and Resources statement.
% Include in any problem set with:  % Shared AI and Resources statement.
% Include in any problem set with:  % Shared AI and Resources statement.
% Include in any problem set with:  \input{ai-statement}
\section*{AI and Resources statement}

Please list (in detail) all resources you used for this assignment. If
you worked with people, list them here as well. It is not enough to say
that you used a resource for help, you need to be specific on the link
and \emph{how} it was helpful. W/R/T gen AI tools (including GPT,
Claude, etc.) you cannot use them to do work on your behalf---you
cannot put in any of the questions, etc. You can ask for help on logic
/ sample problems. If you do use GPT or other AI tools, you need to
provide a link to your chat transcript. Any suspected academic
integrity violations will be immediately reported to the Dean of
Students.

\section*{AI and Resources statement}

Please list (in detail) all resources you used for this assignment. If
you worked with people, list them here as well. It is not enough to say
that you used a resource for help, you need to be specific on the link
and \emph{how} it was helpful. W/R/T gen AI tools (including GPT,
Claude, etc.) you cannot use them to do work on your behalf---you
cannot put in any of the questions, etc. You can ask for help on logic
/ sample problems. If you do use GPT or other AI tools, you need to
provide a link to your chat transcript. Any suspected academic
integrity violations will be immediately reported to the Dean of
Students.

\section*{AI and Resources statement}

Please list (in detail) all resources you used for this assignment. If
you worked with people, list them here as well. It is not enough to say
that you used a resource for help, you need to be specific on the link
and \emph{how} it was helpful. W/R/T gen AI tools (including GPT,
Claude, etc.) you cannot use them to do work on your behalf---you
cannot put in any of the questions, etc. You can ask for help on logic
/ sample problems. If you do use GPT or other AI tools, you need to
provide a link to your chat transcript. Any suspected academic
integrity violations will be immediately reported to the Dean of
Students.


\section*{Survey questions}

\begin{enumerate}[label=\arabic*.]
\item How prepared do you feel for the final exam? not prepared 1 2 3
  4 5 very prepared
\item Which content area of the course was the easiest for you? Why?
  Please be specific.
\item Which content area of the course was the most challenging for
  you? Why? Please be specific.
\end{enumerate}

\end{document}